% =====================================================================
%  Intro to LaTeX: Übungsblatt-Vorlage (Aufgabe 1 aus dem Praxisteil)
%  Kompilieren mit pdfLaTeX (Overleaf-Standard)
%  Alles mit "HIER ANPASSEN" an euch anpassen, den Rest einfach übernehmen.
% =====================================================================
\documentclass[a4paper,11pt]{article}

% ---------------------------------------------------------------------
%  Pakete
% ---------------------------------------------------------------------
\usepackage[T1]{fontenc}
\usepackage[ngerman]{babel}        % deutsche Silbentrennung und Bezeichnungen
\usepackage[a4paper, margin=2.5cm, headheight=14pt]{geometry}
\usepackage{amsmath,amssymb}       % align, \text, Symbole
\usepackage[locale=DE]{siunitx}    % \qty{9.81}{\metre\per\second\squared}
\usepackage{fancyhdr}              % Bonus: Kopf- und Fußzeile

% ---------------------------------------------------------------------
%  Eure Angaben (HIER ANPASSEN)
% ---------------------------------------------------------------------
\newcommand{\name}{Max Mustermann}
\newcommand{\matrikel}{1234567}
\newcommand{\gruppe}{Gruppe 3}
\newcommand{\vorlesung}{Experimentalphysik I}
\newcommand{\blatt}{1}

% ---------------------------------------------------------------------
%  Bonus: eigene Makros
% ---------------------------------------------------------------------
\newcommand{\diff}{\mathrm{d}}                    % aufrechtes d für Ableitungen
\newcommand{\abl}[2]{\frac{\diff #1}{\diff #2}}   % \abl{v}{t} ergibt dv/dt

% ---------------------------------------------------------------------
%  Bonus: Kopf- und Fußzeile
% ---------------------------------------------------------------------
\pagestyle{fancy}
\fancyhf{}                                        % alte Einträge löschen
\lhead{\vorlesung{} -- Blatt \blatt}
\rhead{\name}
\cfoot{\thepage}

% ---------------------------------------------------------------------
%  Titel
% ---------------------------------------------------------------------
\title{\vorlesung\\[0.3em] \large Übungsblatt \blatt}
\author{\name \\ Matrikelnummer: \matrikel \\ Übungsgruppe: \gruppe}
\date{\today}

\begin{document}
\maketitle
\thispagestyle{fancy}

% =====================================================================
\section*{Aufgabe 1}
% =====================================================================
Ein Stein fällt aus der Höhe $h = \qty{20}{\metre}$ aus der Ruhe zu Boden.
Der Luftwiderstand wird vernachlässigt.

\subsection*{a)}
Die Bewegungsgleichung lautet
\begin{equation}
  \abl{v}{t} = -g
  \label{eq:bewegung}
\end{equation}
mit der Erdbeschleunigung $g = \qty{9.81}{\metre\per\second\squared}$.

\subsection*{b)}
Zweimaliges Integrieren von Gleichung~\eqref{eq:bewegung} liefert
$z(t) = h - \frac{1}{2} g t^2$. Für die Fallzeit gilt $z(t_\text{F}) = 0$:
\begin{align}
  0 &= h - \frac{1}{2} g t_\text{F}^2 \\
  \Rightarrow \quad t_\text{F} &= \sqrt{\frac{2h}{g}}
    = \sqrt{\frac{2 \cdot \qty{20}{\metre}}{\qty{9.81}{\metre\per\second\squared}}}
    \approx \qty{2.02}{\second}
\end{align}

\subsection*{c)}
Die Geschwindigkeit beim Aufprall ist
\begin{align*}
  v_\text{F} &= g \, t_\text{F} = \sqrt{2 g h} \\
             &\approx \qty{19.8}{\metre\per\second}.
\end{align*}

% =====================================================================
\section*{Aufgabe 2}
% =====================================================================
\subsection*{a)}
% Hier geht es weiter ...

\end{document}
